\documentclass[11pt,a4paper]{article}

\usepackage[T1]{fontenc}
\usepackage[utf8]{inputenc}
\usepackage[ngerman]{babel}

\usepackage{lmodern}
\usepackage{microtype}
\usepackage{geometry}
\geometry{margin=2.7cm}

\usepackage{amsmath,amssymb,amsthm,mathtools}
\usepackage{bm}
\usepackage{enumitem}
\usepackage{booktabs}
\usepackage{array}
\usepackage{xcolor}
\usepackage{graphicx}
\DeclareGraphicsExtensions{.pdf,.png,.jpg}
\graphicspath{{./}{fig/}{figs/}{images/}{img/}}
\newcommand{\includegraphicssafe}[2][]{\IfFileExists{#2}{\includegraphics[#1]{#2}}{\fbox{\texttt{Missing graphic: #2}}}}

\usepackage[hidelinks]{hyperref}
\hypersetup{
  pdftitle={Ptolemäische, Walter- und Morley-Strukturen auf Primzahlvierlingen und ihrer Dehnungsfamilie},
  pdfauthor={Thomas Hoffbauer},
  pdfsubject={Viererordnungen, Primzahlvierlinge, Ptolemäus, Morley, Marion Walter},
  pdfkeywords={Primzahlvierlinge, Ptolemäus, Morley, Marion Walter, Dehnungsfamilie, Bamberg-Modell}
}

\newtheorem{satz}{Satz}
\newtheorem{lemma}[satz]{Lemma}
\newtheorem{korollar}[satz]{Korollar}
\theoremstyle{definition}
\newtheorem{definition}[satz]{Definition}
\newtheorem{beispiel}[satz]{Beispiel}
\theoremstyle{remark}
\newtheorem{bemerkung}[satz]{Bemerkung}

\newcommand{\R}{\mathbb{R}}
\newcommand{\N}{\mathbb{N}}
\newcommand{\dd}{\mathrm{d}}
\newcommand{\Area}{\mathrm{Area}}
\newcommand{\Mor}{\mathrm{Mor}}
\newcommand{\Wal}{\mathrm{Wal}}

\title{\textbf{Ptolemäische, Walter- und Morley-Strukturen auf Primzahlvierlingen\\ und ihrer Dehnungsfamilie}\\[0.4em]
\large Ein Standalone-Artikel zur geometrischen Klassifikation kompatibler Viererordnungen}
\author{Thomas Hoffbauer}
\date{\today}

\begin{document}

\maketitle

\begin{abstract}
Wir betrachten Primzahlvierlinge der Form
\[
Q_0(p)=(p,p+2,p+6,p+8)
\]
als \emph{primale Grundform} einer Viererordnung mit zwei Zwillingsflügeln und einer festen mod-$12$-Klassenarchitektur. Ausgehend von einer logarithmischen Einbettung sowie einer EABC-Einbettung in die Ebene werden drei komplementäre Klassifikationsinstrumente eingeführt: der ptolemäische Defekt als Maß globaler Rahmenverträglichkeit, Walter-Splits als Maß dehnungsinduzierter Besetzungsasymmetrie und Morley-Splits als Maß lokaler Gestaltreife kanonischer Triangulierungen. Die Analyse wird auf die Dehnungsfamilie
\[
Q_k(p)=(p,p+2,p+6+12k,p+8+12k), \qquad k\in\N_0,
\]
erweitert. Für die lineare Kreuzform wird ein explizites asymptotisches Skalengesetz bewiesen:
\[
\Delta E_{\mathrm{metrisch}}(p,k)\sim \frac{64+288k+288k^2}{p^2}.
\]
Ferner werden explizite asymptotische Walter-Gesetze hergeleitet, während die Morley-Profile als asymptotisch heilende Gestaltkorrekturen interpretiert werden. Der Artikel ordnet die Entdeckung als eine hierarchische Familie strukturtreuer Deformationen derselben Viererarchitektur ein.
\end{abstract}

\section{Motivation}

Primzahlvierlinge
\[
(p,p+2,p+6,p+8)
\]
sind die dichtesten möglichen Viererkonfigurationen ungerader Primzahlen. Arithmetisch werden sie meist unter dem Gesichtspunkt der Primzahllücken oder der Hardy--Littlewood-Heuristik betrachtet. Die leitende Idee dieses Artikels ist eine andere: Ein Primzahlvierling wird hier nicht nur als Menge von vier Zahlen verstanden, sondern als \emph{geometrisch strukturierte Viererordnung}.

Der Ausgangspunkt ist die Beobachtung, dass ein Primzahlvierling zugleich drei Eigenschaften besitzt:
\begin{enumerate}[label=(\arabic*)]
\item zwei \emph{Zwillingsflügel} $(p,p+2)$ und $(p+6,p+8)$,
\item eine natürliche lineare Größenordnung,
\item eine stabile Klassenarchitektur modulo $12$, die in der Sprache dieses Artikels als
\[
A,B,C,E
\]
bezeichnet wird.
\end{enumerate}

Diese drei Ebenen erlauben es, geometrische Invarianten und Defekte einzuführen. Insbesondere entstehen drei Fragen:
\begin{itemize}
\item Ist eine Viererordnung \emph{ptolemäisch kompatibel}, also rahmenverträglich?
\item Wie verteilt sich ihre \emph{Besetzung} auf natürliche Triangulierungen?
\item Wie reif ist ihre lokale \emph{Gestalt} im Sinn von Morley?
\end{itemize}

Aus diesen Fragen entsteht eine neue Sichtweise: Der klassische Primzahlvierling erscheint als \emph{Urform} einer größeren Dehnungsfamilie mit erhaltener Klassenarchitektur,
\[
Q_k(p)=(p,p+2,p+6+12k,p+8+12k).
\]
Der Dehnungsrang $k$ verändert die Korrekturterme, nicht aber die zugrunde liegende Viererarchitektur.

\section{Die primale Grundform und ihre Dehnungsfamilie}

\begin{definition}[Primale Grundform]
Ein Primzahlvierling der Form
\[
Q_0(p):=(p,p+2,p+6,p+8)
\]
heißt \emph{primale Grundform}, sofern alle vier Einträge prim sind.
\end{definition}

\begin{definition}[Dehnungsfamilie]
Für $k\in\N_0$ definieren wir
\[
m_k:=6+12k
\]
und
\[
Q_k(p):=(p,p+2,p+m_k,p+m_k+2).
\]
Sofern alle vier Einträge prim sind, heißt $Q_k(p)$ eine \emph{Viererordnung vom Dehnungsrang $k$}.
\end{definition}

Die Bedingung $m_k\equiv 6\pmod{12}$ stellt sicher, dass die mod-$12$-Struktur der Grundform erhalten bleibt. Die ersten Stufen sind:
\[
Q_0(p)=(p,p+2,p+6,p+8),
\]
\[
Q_1(p)=(p,p+2,p+18,p+20),
\]
\[
Q_2(p)=(p,p+2,p+30,p+32),
\]
\[
Q_3(p)=(p,p+2,p+42,p+44).
\]

\begin{bemerkung}
Im Bamberg-Vokabular ist $Q_0$ die \emph{Urstufe} der Viererordnung, während $Q_k$ für $k\ge 1$ strukturtreue \emph{Dehnungsnachfahren} derselben Rahmenarchitektur sind.
\end{bemerkung}

\section{Zwei Geometrien auf derselben Viererordnung}

\subsection{Lineare logarithmische Geometrie}

Wir definieren
\[
x_1=\log p,\qquad x_2=\log(p+2),\qquad x_3=\log(p+m_k),\qquad x_4=\log(p+m_k+2).
\]
Diese vier Punkte liegen in natürlicher Reihenfolge auf der reellen Achse.

\subsection{EABC-Einbettung}

Wir kodieren die Klassen $E,A,B,C$ als vier orthogonale Richtungen:
\[
E\mapsto (1,0),\qquad A\mapsto (0,1),\qquad B\mapsto (-1,0),\qquad C\mapsto (0,-1).
\]
Mit logarithmischen Radien erhalten wir die Punkte
\[
A_p=(0,\log p),\qquad B_p=(-\log(p+2),0),
\]
\[
C_p=(0,-\log(p+m_k)),\qquad E_p=(\log(p+m_k+2),0).
\]

\begin{definition}[Kompatible Form und Kreuzform]
Zur linearen oder EABC-Geometrie gehören zwei Anordnungen:
\[
\text{kompatible Form}: (1,2,3,4)\ \text{bzw.}\ (A,B,C,E),
\]
\[
\text{Kreuzform}: (1,3,2,4)\ \text{bzw.}\ (A,C,B,E).
\]
\end{definition}

\section{Ptolemäische Strukturgesetze}

\begin{definition}[Ptolemäischer Defekt]
Für vier Punkte $X_1,X_2,X_3,X_4\in \R^n$ definieren wir
\[
\Pi(X_1,X_2,X_3,X_4)
=
d_{13}d_{24}-d_{12}d_{34}-d_{23}d_{41},
\]
wobei $d_{ij}=\|X_i-X_j\|$.
Der Betrag $|\Pi|$ misst die Abweichung von ptolemäischer Balance.
\end{definition}

\subsection{Exakter linearer Satz}

\begin{satz}[Exakte lineare Schließung der kompatiblen Form]
Seien $x_1<x_2<x_3<x_4$ Punkte auf der reellen Achse. Dann gilt
\[
\Pi(x_1,x_2,x_3,x_4)=0.
\]
Insbesondere gilt für jede Viererordnung $Q_k(p)$:
\[
\Pi_{\mathrm I}^{\mathrm{lin}}(p,k)=0.
\]
\end{satz}

\begin{proof}
Auf der Geraden sind
\[
d_{13}=x_3-x_1,\quad d_{24}=x_4-x_2,\quad d_{12}=x_2-x_1,\quad d_{34}=x_4-x_3,
\]
\[
d_{23}=x_3-x_2,\quad d_{41}=x_4-x_1.
\]
Daher
\begin{align*}
\Pi
&=(x_3-x_1)(x_4-x_2)-(x_2-x_1)(x_4-x_3)-(x_3-x_2)(x_4-x_1)\\
&=0
\end{align*}
nach direkter Ausmultiplizierung und Kürzung aller Terme.
\end{proof}

\begin{satz}[Expliziter linearer Kreuzdefekt]
Für die Kreuzform $(x_1,x_3,x_2,x_4)$ gilt
\[
\Pi_{\mathrm{II}}^{\mathrm{lin}}(p,k)
=
-2\log\!\Bigl(\frac{p+m_k}{p+2}\Bigr)\log\!\Bigl(\frac{p+m_k+2}{p}\Bigr).
\]
Insbesondere gilt für festes $k$ und $p\to\infty$:
\[
\Delta E_{\mathrm{metrisch}}(p,k):=
\bigl|\Pi_{\mathrm{II}}^{\mathrm{lin}}(p,k)\bigr|
\sim
\frac{64+288k+288k^2}{p^2}.
\]
\end{satz}

\begin{proof}
Aus der Definition der Kreuzform mit $x_1=\log p$, $x_2=\log(p+2)$, $x_3=\log(p+m_k)$, $x_4=\log(p+m_k+2)$ folgt direkt
\[
\Pi_{\mathrm{II}}^{\mathrm{lin}}(p,k)
=
-2(x_3-x_2)(x_4-x_1),
\]
also die angegebene Formel. Für die Asymptotik benutzen wir
\[
\log\!\Bigl(\frac{p+m_k}{p+2}\Bigr)\sim \frac{m_k-2}{p},
\qquad
\log\!\Bigl(\frac{p+m_k+2}{p}\Bigr)\sim \frac{m_k+2}{p},
\]
also
\[
\Delta E_{\mathrm{metrisch}}(p,k)\sim \frac{2(m_k^2-4)}{p^2}.
\]
Mit $m_k=6+12k$ ist
\[
2(m_k^2-4)=2\bigl((6+12k)^2-4\bigr)=64+288k+288k^2.
\]
\end{proof}

\begin{bemerkung}
Der lineare Satz trennt zwei Ebenen sauber: Die kompatible Form ist metrisch exakt geschlossen; die Kreuzform erzeugt einen expliziten Defekt, der quadratisch in $k$ wächst und zugleich mit $p^{-2}$ verschwindet.
\end{bemerkung}

\subsection{EABC-Strukturspannung}

\begin{definition}[EABC-Strukturspannung]
Für die EABC-Kreuzform definieren wir
\[
\Delta E_{EABC}(p,k):=
\bigl|\Pi(A_p,C_p,B_p,E_p)\bigr|.
\]
\end{definition}

\begin{satz}[Universeller Hauptterm der EABC-Kreuzspannung]
Für festes $k$ und $p\to\infty$ gilt formal-asymptotisch
\[
\Delta E_{EABC}(p,k)\sim 4(\log p)^2.
\]
\end{satz}

\begin{proof}[Beweisidee]
Für festes $k$ werden die vier Radien
\[
\log p,\ \log(p+2),\ \log(p+m_k),\ \log(p+m_k+2)
\]
asymptotisch gleich. Die EABC-Figur nähert sich daher einem symmetrischen Kreuzviereck mit Radius $\log p$. In dieser Grenzfigur trägt die Kreuzverschaltung $(A,C,B,E)$ einen führenden Defektterm der Größenordnung $4(\log p)^2$. Die $k$-Abhängigkeit erscheint nur in niedrigeren Korrekturtermen, weil sie lediglich Radialdifferenzen der Größenordnung $O(1/p)$ erzeugt.
\end{proof}

\begin{bemerkung}
Hier liegt der konzeptionell wichtigste Unterschied zur linearen Geometrie: Der metrische Kreuzdefekt ist ein endlicher Dehnungsrest, die EABC-Strukturspannung dagegen ein struktureller Rahmeneffekt.
\end{bemerkung}

\section{Walter-Kanäle als Besetzungsasymmetrie}

Die kompatible EABC-Figur besitzt zwei natürliche Triangulierungen:
\[
\triangle ABC,\qquad \triangle ACE,
\]
\[
\triangle ABE,\qquad \triangle BCE.
\]

\begin{definition}[Walter-Wert]
Sei $T$ ein Dreieck. Der Walter-Wert $\Wal(T)$ sei proportional zur Fläche des zentralen Walter-Hexagons. Da Marion Walters Satz einen festen Flächenanteil liefert, genügt es in diesem Artikel, $\Wal(T)$ als proportional zu $\Area(T)$ zu behandeln.
\end{definition}

\begin{definition}[Walter-Splits]
Wir setzen
\[
\Delta W_{AC}(p,k):=\Wal(\triangle ACE)-\Wal(\triangle ABC),
\]
\[
\Delta W_{BE}(p,k):=\Wal(\triangle BCE)-\Wal(\triangle ABE).
\]
\end{definition}

\begin{satz}[Explizite Walter-Formeln]
Für die Dehnungsfamilie gilt
\[
\Delta W_{AC}(p,k)
=
\frac{1}{20}\bigl(\log p+\log(p+m_k)\bigr)\bigl(\log(p+m_k+2)-\log(p+2)\bigr),
\]
\[
\Delta W_{BE}(p,k)
=
\frac{1}{20}\bigl(\log(p+2)+\log(p+m_k+2)\bigr)\bigl(\log(p+m_k)-\log p\bigr).
\]
Insbesondere gilt für festes $k$ und $p\to\infty$:
\[
\Delta W_{AC}(p,k)\sim \Delta W_{BE}(p,k)\sim \frac{(6+12k)\log p}{10p}.
\]
\end{satz}

\begin{proof}
Die Flächen der beteiligten Dreiecke ergeben sich direkt aus den Koordinaten in der EABC-Einbettung. Wegen der linearen Proportionalität des Walter-Werts zur Dreiecksfläche genügt eine elementare Flächenberechnung. Für die Asymptotik verwenden wir
\[
\log(p+m_k+2)-\log(p+2)\sim \frac{m_k}{p},
\qquad
\log(p+m_k)-\log p\sim \frac{m_k}{p},
\]
sowie
\[
\log p+\log(p+m_k)\sim 2\log p,\qquad
\log(p+2)+\log(p+m_k+2)\sim 2\log p.
\]
Daraus folgt
\[
\Delta W_{AC}(p,k)\sim \Delta W_{BE}(p,k)\sim \frac{m_k\log p}{10p}.
\]
Mit $m_k=6+12k$ ergibt sich die Behauptung.
\end{proof}

\begin{bemerkung}
Walter misst im vorliegenden Modell nicht grobe Fehlverschaltung, sondern die \emph{Besetzungsasymmetrie} der Dehnung. Der Split wächst linear mit $k$.
\end{bemerkung}

\section{Morley-Profile als Gestaltreife}

\begin{definition}[Morley-Quotient]
Für ein Dreieck $T$ definieren wir
\[
M(T):=\frac{\Area(\Mor(T))}{\Area(T)},
\]
wobei $\Mor(T)$ das Morley-Dreieck von $T$ bezeichnet.
\end{definition}

\begin{definition}[Morley-Splits]
Für die kanonischen Triangulierungen setzen wir
\[
\Delta M_{AC}(p,k):=M(\triangle ACE)-M(\triangle ABC),
\]
\[
\Delta M_{BE}(p,k):=M(\triangle BCE)-M(\triangle ABE).
\]
\end{definition}

\begin{satz}[Asymptotische Morley-Heilung]
Für festes $k$ und $p\to\infty$ gilt
\[
\Delta M_{AC}(p,k)\to 0,\qquad \Delta M_{BE}(p,k)\to 0.
\]
\end{satz}

\begin{proof}
Für festes $k$ werden die vier logarithmischen Radien asymptotisch gleich. Daher konvergieren alle kanonischen Dreiecke gegen dieselbe asymptotische Referenzgestalt eines gleichschenklig-rechtwinkligen Dreiecks innerhalb der EABC-Einbettung. Da der Morley-Quotient stetig von den Winkeln des Dreiecks abhängt, konvergieren alle vier Morley-Quotienten gegen denselben Grenzwert. Folglich verschwinden ihre Differenzen.
\end{proof}

\begin{bemerkung}
Morley misst lokale Formspannung, nicht globale Rahmenverspannung. Während Walter einen Massenkanal beschreibt, codiert Morley die \emph{Gestaltreife} der Triangulierungen.
\end{bemerkung}

\section{Hierarchie der ersten Dehnungsstufen}

Die ersten Stufen der Dehnungsfamilie lauten:
\[
Q_0(p)=(p,p+2,p+6,p+8),
\]
\[
Q_1(p)=(p,p+2,p+18,p+20),
\]
\[
Q_2(p)=(p,p+2,p+30,p+32),
\]
\[
Q_3(p)=(p,p+2,p+42,p+44),
\]
\[
Q_4(p)=(p,p+2,p+54,p+56).
\]

Dazu gehören die asymptotischen linearen Kreuzreste
\[
\Delta E_{\mathrm{metrisch}}(p,0)\sim \frac{64}{p^2},
\qquad
\Delta E_{\mathrm{metrisch}}(p,1)\sim \frac{640}{p^2},
\]
\[
\Delta E_{\mathrm{metrisch}}(p,2)\sim \frac{1792}{p^2},
\qquad
\Delta E_{\mathrm{metrisch}}(p,3)\sim \frac{3520}{p^2},
\qquad
\Delta E_{\mathrm{metrisch}}(p,4)\sim \frac{5824}{p^2}.
\]

Und für die Walter-Kanäle gilt
\[
\Delta W(p,0)\sim \frac{3\log p}{5p},
\qquad
\Delta W(p,1)\sim \frac{9\log p}{5p},
\]
\[
\Delta W(p,2)\sim \frac{3\log p}{p},
\qquad
\Delta W(p,3)\sim \frac{21\log p}{5p},
\qquad
\Delta W(p,4)\sim \frac{27\log p}{5p}.
\]

\begin{center}
\begin{tabular}{@{}>{$}c<{$} >{$}c<{$} >{$}c<{$} >{$}c<{$}@{}}
\toprule
k & m_k & Q_k(p) & \Delta E_{\mathrm{metrisch}}(p,k) \\
\midrule
0 & 6  & (p,p+2,p+6,p+8)   & 64/p^2 \\
1 & 18 & (p,p+2,p+18,p+20) & 640/p^2 \\
2 & 30 & (p,p+2,p+30,p+32) & 1792/p^2 \\
3 & 42 & (p,p+2,p+42,p+44) & 3520/p^2 \\
4 & 54 & (p,p+2,p+54,p+56) & 5824/p^2 \\
\bottomrule
\end{tabular}
\end{center}

\section{Klassifikationsvektor}

Die drei Ebenen --- Ptolemäus, Walter, Morley --- lassen sich zu einem Profil bündeln.

\begin{definition}[Klassifikationsvektor]
Der Klassifikationsvektor der Viererordnung $Q_k(p)$ sei
\[
\mathcal K(Q_k(p))
=
\bigl(
k,\,
\Delta E_{\mathrm{metrisch}}(p,k),\,
\Delta E_{EABC}(p,k),\,
\Delta W_{AC}(p,k),\,
\Delta W_{BE}(p,k),\,
\Delta M_{AC}(p,k),\,
\Delta M_{BE}(p,k)
\bigr).
\]
\end{definition}

\begin{bemerkung}
Dieser Vektor trennt:
\begin{itemize}
\item den \emph{Dehnungsrang} $k$,
\item den \emph{metrischen Rest} der linearen Kreuzform,
\item die \emph{klassengeometrische Rahmenverspannung},
\item die \emph{Besetzungsasymmetrie} der Walter-Kanäle,
\item die \emph{lokale Gestaltreife} der Morley-Profile.
\end{itemize}
\end{bemerkung}

\section{Einordnung der Entdeckung}

Die hier vorgestellte Struktur besitzt aus Sicht des Autors drei neuartige Züge.

\subsection*{(i) Vom Objekt zur Familie}
Der klassische Primzahlvierling wird nicht isoliert betrachtet, sondern als erste Stufe einer unendlichen Dehnungsfamilie
\[
Q_k(p)=(p,p+2,p+6+12k,p+8+12k).
\]
Damit erscheint $m=6$ nicht mehr als singulärer Sonderfall, sondern als \emph{Urform}.

\subsection*{(ii) Dreifache Geometrisierung}
Drei klassische geometrische Themen werden auf dieselbe arithmetische Struktur gelegt:
\begin{itemize}
\item \textbf{Ptolemäus} als Test globaler Rahmenverträglichkeit,
\item \textbf{Walter} als Sensor der Besetzungsasymmetrie,
\item \textbf{Morley} als Maß lokaler Gestaltreife.
\end{itemize}
Diese Trias erzeugt eine Klassifikation, die über die bloße Angabe von Primzahllücken hinausgeht.

\subsection*{(iii) Trennung von Rahmen und Korrektur}
Die Untersuchung legt nahe:
\begin{itemize}
\item Die \emph{Rahmenstruktur} bleibt über alle Dehnungsstufen hinweg invariant.
\item Die \emph{Korrekturen} --- metrisch, besetzungsbezogen, lokalgestaltlich --- wachsen oder verschwinden nach klaren asymptotischen Gesetzen.
\end{itemize}
Gerade diese Trennung von invariantem Rahmen und variabler Korrektur ist der eigentliche konzeptionelle Kern der Entdeckung.

\begin{bemerkung}
Die Aussagen zu Walter und insbesondere zu Morley sind als strukturtheoretische und asymptotische Aussagen formuliert. Für allgemeine $k$ wäre eine breite numerische Auswertung eigenständiges Folgematerial. Der lineare Ptolemäus-Satz und die Walter-Asymptotik sind hingegen bereits in expliziter Form vorhanden.
\end{bemerkung}

\section{Diskussion}

Die Dehnungsfamilie legt eine Sicht nahe, in der Primzahlvierlinge nicht nur arithmetische Seltenheiten sind, sondern Träger einer internen Viererarchitektur. In dieser Sprache bedeutet:
\begin{itemize}
\item $Q_0$ ist die \emph{kanonische Urform},
\item $Q_k$ für $k\ge 1$ sind \emph{strukturtreue Dehnungsnachfahren},
\item die kompatible Form ist der bevorzugte Rahmen,
\item die Kreuzform misst Verspannung gegen diesen Rahmen.
\end{itemize}

Die drei Klassifikationsinstrumente greifen dabei auf unterschiedliche Ebenen:
\begin{center}
\begin{tabular}{@{}ll@{}}
\toprule
Instrument & Funktion \\
\midrule
Ptolemäus & globale Rahmenverträglichkeit \\
Walter & Besetzungs- und Lastverteilung \\
Morley & lokale Gestalt- und Winkelreife \\
\bottomrule
\end{tabular}
\end{center}

\section{Ausblick}

Es bieten sich mindestens vier nächste Schritte an:
\begin{enumerate}[label=(\arabic*)]
\item eine numerische Vollauswertung der Familien $Q_k$ für kleine $k$,
\item eine explizite Bestimmung der Unterterme in $\Delta E_{EABC}(p,k)$,
\item eine präzise numerische Skalierung der Morley-Profile,
\item die Einbettung dieser Familie in ein allgemeineres Vierer- oder Mehrflügelmodell.
\end{enumerate}

\section*{Schlussformel}

Die Primzahlvierlinge
\[
Q_0(p)=(p,p+2,p+6,p+8)
\]
bilden die primale Grundform einer unendlichen Familie strukturtreuer Dehnungen
\[
Q_k(p)=(p,p+2,p+6+12k,p+8+12k).
\]
Über alle Stufen bleibt die Viererarchitektur erhalten. Variabel sind nur die endlichen Korrekturterme: Der metrische Kreuzrest wächst quadratisch mit dem Dehnungsrang, die Walter-Besetzung linear, während die Morley-Verzerrung asymptotisch heilt. Die Entdeckung besteht damit nicht nur in einer einzelnen Formel, sondern in der Sichtbarkeit einer ganzen hierarchischen Geometrie von Viererordnungen.

\end{document}
